\section{Adoption order and bounded follow-up work}

\subsection{Recommended first pass}

The coefficient patch is the strongest immediate candidate. It removes
repeated mathematical expansion, applies under weaker local hypotheses,
serves five existing proof sites, and saves 20 physical lines after the
shared theorem's full cost. Its new module is small and has an explicit
interface. Review the diagonal multiplicity and the selected exponent's
fibre condition first; they are the essential correctness points.

The Pell patch is the smallest isolated proof edit. It adds neither an
import nor a helper and saves eight lines. Its mathematical obligation is
unchanged: the same interval between consecutive squares rules out the
small second coordinate. The change is to derive that coordinate's square
once from its absolute value. It is a suitable independent integration
commit.

The codec/parent patch is useful where preserving the literal wire format
is a priority. It makes an exact relationship to a library API explicit and
turns an already available raw correctness theorem into the source of the
typed correctness theorem. The four-line net saving is smaller, but the
duplicate induction disappears. The optional length proof can be considered
separately after examining its added \code{Classical.choice} dependency.

Defer the first-failure patch if immediate aggregate source length is the
only criterion. Adopt it when shared control of interval minimality is
valuable, or when an additional strip argument would otherwise introduce
another copy. No arbitrary threshold on a predicted future saving is needed:
compare the actual next client against the 28-line helper already supplied.

\subsection{Two checked prototypes, with no unearned production saving}

\textbf{Mixed products.} The optional \code{PairFiber.lean} file supplies a
related local coefficient theorem for a product of two encodings. At the
chosen exponent the two surviving contributions are \(u_i v_j\) and
\(u_j v_i\); on the diagonal there is only \(u_i v_i\). It can support
future bilinear clients that should not be forced through a square. This
prototype was compiled, but no present production client is migrated to
it and no savings are counted for it. Integrate only when a concrete client
justifies the extra interface.

\textbf{Pell minimum-coordinate criterion.} The optional
\code{PellMinimum.lean} file proves that minimizing \(y^2\) among solutions
with \(x>1\), together with the candidate's positivity, suffices for
fundamentality. It makes the mathematical proof obligation reusable while
allowing competitors of either sign in \(y\). For the one current target,
the direct eight-line reduction is preferable. A second fundamental-unit
proof would give a concrete reason to revisit the shared criterion.

Both prototypes are accompanied by successful compiler logs. They are
separate from the ten-module production overlay and from all source-saving
totals. Neither is a reason to reopen an already-excluded historical idea.

\subsection{What would change these recommendations}

Three concrete developments could affect the choice. A changed Mathlib
list-encoding instance would require rechecking the exact codec bridge.
Another polynomial encoding with only selected isolated targets would
increase the value of the local coefficient theorem. Another strip-wall
proof would provide a real third client for the first-failure abstraction.
Each can be assessed through the supplied interfaces without redesigning
the executable reductions.

The archive is meant to make those decisions reviewable: each direction has
an exact target, a mathematical reason for the simplification, a separate
integration patch, complete affected source, and a declared verification
scope. The exclusion ledger accompanies the proposals so that later work can
continue from this boundary instead of rediscovering the same mechanisms.
